\subsubsection{Local theorem dependencies}\label{ssec:theorem-lanes}
This section records local hypotheses, asserted scope, and
proof inputs.  The proof-bearing text is the cited theorem,
proposition, lemma, or computation.  Each assertion is interpreted with
the proof status and admissibility hypotheses recorded in the claim-status table
and the stratified Koszul criterion.

\begin{stmt}[Statement: Dirac brane probe]\label{thm:lane-dirac-probe}
  \emph{Status.} \emph{proved in finite algebraic model}.  The proof is
  confined to the formal local ghost-zero sector of the brane probe.

  \emph{Exact hypotheses.} The transverse formal disk is
  \((\widehat{\C^2}_0,dz_1\wedge dz_2)\); the stack size \(N\) is finite;
  fields are \(\phi_1,\phi_2,A\in\mathfrak{gl}_N\) on the brane line;
  the raw shifted-Chern--Simons coefficient fields have first been put in
  the Dirac normalization of
  Construction~\ref{constr:open-field-basis}; variations are compactly
  supported; and the only residual gauge group retained in this statement is
  ordinary \(GL_N\)-conjugation on degree-zero functions.

  \emph{Local assertion.} For a stack of \(N\) branes probing
  \((\widehat{\C^2}_0,dz_1\wedge dz_2)\), the ghost-zero open action is
  the first-order Dirac matrix system
  \[
    S_\partial
    =\int_\R\operatorname{Tr}
      \bigl(\phi_1\,d\phi_2+A[\phi_1,\phi_2]\bigr),
  \]
  and its reduced BV algebra is
  \[
    \C[\phi_1{}^i{}_j,\phi_2{}^i{}_j]\otimes
    \Lambda(\psi^i{}_j),
    \qquad
    Q\psi=[\phi_1,\phi_2],\quad Q\phi_i=0.
  \]
  The moment map for conjugation is
  \(-\operatorname{Tr}(\epsilon[\phi_1,\phi_2])\).  With the paper's
  convention \(\iota_{X_f}\Omega=-df\), its components satisfy
  \(\{\mu_\epsilon,\mu_\eta\}=-\mu_{[\epsilon,\eta]}\); this
  anti-equivariance affects only the comoment convention, not the zero
  locus.  The BV differential is the derived Koszul zero-fibre
  differential for the commuting-pair equation, not an underived regular
  sequence assertion.

  \emph{Proof inputs.}
  Lemma~\ref{lem:dirac-probe-reduction} derives the ghost-zero action
  from the shifted normal directions and fixes the moment-map sign.
  Proposition~\ref{prop:open-bv-truncation} proves that the ghost-zero
  truncation is consistent and identifies the residual Koszul complex.

  \emph{Scope and exclusions.} The statement does not assert an untruncated open-string
  BV model beyond this ghost-zero finite algebraic truncation.
\end{stmt}

\begin{stmt}[Statement: unital cyclic Koszul and normal-ordering core]\label{thm:lane-cyclic-koszul-normal-ordering}
  \emph{Status.} \emph{proved algebraically, with an explicit
  obstruction to the full cyclic contraction}.  This is the
  necklace-level Koszul calculation with the scalar word and
  one-\(\psi\) homology kept explicit.

  \emph{Exact hypotheses.} Work in the unital cyclic dg Lie algebra
  \[
    \mathfrak p_{\mathrm K}
    =
    \operatorname{Cyc}\bigl(T(\C x\oplus\C y\oplus\C\psi)\bigr),
    \qquad
    Q\psi=xy-yx,\quad Qx=Qy=0,
  \]
  with the scalar word \([1]\) retained until the bracket has been
  checked.  The scalar-reduced quotient
  \(\bar{\mathfrak p}_{\mathrm K}=\mathfrak p_{\mathrm K}/\C[1]\) is
  taken only after this unital computation.  The normal-ordering
  homotopy requires a chosen cyclic cut and applies to the no-\(\psi\)
  subcomplex selected by that cut.

  \emph{Local assertion.} The cyclic necklace bracket makes
  \(\mathfrak p_{\mathrm K}\) a dg Lie algebra.  The normal-ordering map
  \(N\) is a strict dg Lie map to the unital polynomial Hamiltonian
  algebra and sends \([1]\) to \(1\); after quotienting constants it
  gives the reduced Hamiltonian map
  \(N_{\mathrm{red}}\colon\bar{\mathfrak p}_{\mathrm K}\to\bar A\).
  On the cut no-\(\psi\) sector the straightening homotopy satisfies
  \(QH_c=\mathrm{id}-iN\).  There is no cyclic homotopy
  \(K\) with \(QK+KQ=\mathrm{id}-iN\) on the full cyclic complex when
  \(N\) kills the primitive one-\(\psi\) classes.

  \emph{Proof inputs.}
  Theorem~\ref{thm:primitive-cyclic-koszul-p0-model} proves the unital
  dg Lie model and the scalar bracket \(\{[x],[y]\}=[1]\).
  Proposition~\ref{prop:normal-ordering-quotient-cyclic-koszul}
  proves the strict normal-ordering quotient.
  Theorem~\ref{thm:normal-ordering-straightening-homotopy} constructs
  the cut-dependent no-\(\psi\) straightening homotopy.
  Lemma~\ref{lem:no-full-cyclic-normal-ordering-contraction} proves the
  obstruction to a full cyclic contraction.  Lemma~\ref{lem:gauge-paired-necklace-naive-quotient-fails}
  records the failure of the naive gauge-paired quotient.

  \emph{Scope and exclusions.} This block is not an unreduced
  BV/factorization-centre theorem.  It does not identify arbitrary
  open observables with central operations, and polynomial
  one-\(\psi\) descendants do not replace the principal-part cotangent
  current module.
\end{stmt}

\begin{stmt}[Statement: stable trace and one-\(\psi\) PBW shadow]\label{thm:lane-stable-trace}
  \emph{Status.} \emph{proved degreewise stable}.  It is a PBW
  special-fibre trace statement, not a statement about the principal-part
  cotangent module.

  \emph{Exact hypotheses.} Work at fixed total polynomial degree and then
  in the stable connected large-\(N\) range; pass to scalar-reduced
  connected traces; use the Koszul differential
  \(Q\psi=[\phi_1,\phi_2]\); and take the PBW special fibre before any
  filtered or \(D_\hbar\)-interpolation is introduced.

  \emph{Local assertion.} The Hamiltonian open observables in ghost
  number zero are generated by scalar-reduced connected trace classes
  \(\overline{\operatorname{Tr} f(\phi_1,\phi_2)}\).  The non-scalar
  primitive one-\(\psi\) classes \(\Psi_{a,b}\), \(a+b>0\), are the primitive PBW
  special-fibre shadow classes.  They do not realize the continuous
  \(\mathfrak h\)-coadjoint module and are not the principal-part
  currents \(\rho_{a,b}\).

  \emph{Proof inputs.}
  Procesi--Razmyslov stable invariant theory supplies the stable trace
  generators.  The Koszul differential \(Q\psi=[\phi_1,\phi_2]\) turns
  adjacent swaps into boundaries in the connected trace quotient.
  Proposition~\ref{prop:one-psi-homology} computes the primitive
  one-\(\psi\) homology.  Theorem~\ref{thm:chain-level-primitive-projection}
  identifies the chain-level primitive shadow, while
  Theorem~\ref{thm:pbw-vs-deformation} and
  Theorem~\ref{thm:no-polynomial-realization-categorical} separate that
  PBW shadow from the coadjoint principal-part module.
  Theorem~\ref{thm:one-psi-label-hbar-obstruction} proves the
  same separation for every \(\hbar\)-adic deformation satisfying its
  flatness, PBW-linear-action, Moyal-coadjoint-target, and leading-term
  hypotheses.

  \emph{Scope and exclusions.} Unstable finite-\(N\) trace identities and any
  ordinary \(\mathfrak h\)-module identification of
  \(\Psi_{a,b}\) with \(\rho_{a,b}\) are outside this statement.  The
  Fourier--Rees comparison of
  Theorem~\ref{thm:app-matlis-rees-fourier-bridge} is only an
  associated-graded label bridge for the original Hamiltonian action; the
  legitimate bridge is Fourier-twisted in the Rees
  \(\mathcal D_\hbar\)-module sense.  The corrected quantum target, when
  needed, is the reduced distributional principal-part target of
  Construction~\ref{constr:quantum-principal-part-descendant-target},
  not the PBW one-\(\psi\) span.
\end{stmt}

\begin{stmt}[Statement: CE/PV reduced Hamiltonian central-operation model]\label{thm:lane-ce-pv-center}
  \emph{Status.} \emph{proved degreewise stable}.  The proof is at
  cochain level for the reduced formal Hamiltonian algebra and its
  continuous shifted cotangent extension.  The bracketed \(P_0\)
	  consequence is used only after a bracket-admissible realization has
	  been chosen, and the bar-cobar consequence is used only under the
	  stated lower-central Tate-pronilpotent hypothesis.

  \emph{Exact hypotheses.} The Hamiltonian algebra is
  \(\mathfrak h=\C[[z_1,z_2]]/\C\cdot1\) with its \(m\)-adic topology;
  the cotangent direction is
  \(\mathfrak h^\vee_{\mathrm{cont}}[1]\), equivalently to
  \(\mathcal P[1]\) after the residue scalar is fixed by
  Theorem~\ref{thm:principal-part-coadjoint-uniqueness}; CE cochains are
	  completed in coordinate products; \(P_0\)-brackets require a
	  bracket-admissible subalgebra \(B\); kernels require a
	  kernel-admissible \(B\); and every use of bar-cobar duality is
	  restricted to the lower-central Tate-pronilpotent hypothesis named in
	  the cited lemma.

  \emph{Local assertion.} Let
  \[
    \mathfrak h=\C[[z_1,z_2]]/\C\cdot1,\qquad
    \mathfrak g=\mathfrak h\ltimes
      \mathfrak h^\vee_{\mathrm{cont}}[1].
  \]
  The coordinate completion of the CE cochains of \(\mathfrak g\)
  identifies, windowwise and after the coordinate completion specified in
  Definition~\ref{defn:formal-disk-ce-pv-completion}, with
  the corresponding Schouten polyvector algebra:
  \[
    \Phi_{\mathrm{coord}}\colon
    C^\bullet_{\mathrm{CE},\mathrm{coord}}(\mathfrak g)
    \xrightarrow{\sim}
    \mathrm{PV}_{\mathrm{coord}}
    (A_{\partial,\mathrm{coord}}).
  \]
  On a bracket-admissible \(B\subset
  \mathrm{PV}_{\mathrm{coord}}(A_{\partial,\mathrm{coord}})\) the
  restriction \(\Phi^B_{P_0}\) is a dg \(P_0\)-isomorphism.  On a
  kernel-admissible \(B\) the diagonal tensor \(\Theta_B\) is the
  associated Maurer--Cartan kernel.  The ambient strict product completion
  carries neither assertion.
  The generator dictionary is
  \[
    c^I\mapsto\theta^I,\qquad u_I\mapsto O_I .
  \]
  The \(Z^{P_0}_{\mathrm{red,Ham}}\) identification is asserted only on the
  coordinate-window or mixed/weighted-continuous sector where the Schouten
  bracket is a continuous operation; the strict ambient product
  completion is not a \(P_0\)-algebra here.
  Here \(|u_I|=0\), \(|c^I|=1\), and \(u_I\) is the CE coordinate dual
  to the shifted-cotangent generator, not to a Hamiltonian generator.

  \emph{Proof inputs.}
  Proposition~\ref{prop:grav-ops} identifies closed local operators as
  continuous CE cochains, which sit inside the coordinate envelope only
  after the chosen continuity condition is specified.
  Lemma~\ref{lem:formal-disk-ce-schouten-continuity} computes the
  coordinate PV differential.  Theorem~\ref{thm:reduced-ce-pv-central-operation}
  checks
  \(\Phi_{\mathrm{coord}} d_{\mathrm{CE}}
  =d_\pi^{\mathrm{coord}}\Phi_{\mathrm{coord}}\) on generators and extends
  by the completed Leibniz rule.  Lemma~\ref{lem:continuous-bar-cobar}
  supplies the conilpotent bar-cobar template only in its own
  hypothesis; the untruncated formal Hamiltonian algebra with linear
  modes is not used as a lower-central Tate-pronilpotent input.

  \emph{Scope and exclusions.} The statement does not prove an unrestricted
  factorization-centre theorem or a conilpotence statement for the
  untruncated formal Hamiltonian algebra.  It also does not construct a
  BV heat-kernel propagator, a BV Laplacian, or a QME-compatible
  closed-open graph model for the unreduced cotangent sector.
\end{stmt}

\paragraph{Cochain-level CE/PV--Koszul block.}

\begin{defn}[Cotangent CE/PV datum]\label{defn:cotangent-ce-pv-datum}
  A cotangent CE/PV datum is a Lie algebra \(\mathfrak l\) in a
  complete linear category, a continuous dual module
  \(\mathfrak l^\vee\), and the shifted cotangent Lie algebra
	  \[
	    T^*[1]\mathfrak l=\mathfrak l\ltimes\mathfrak l^\vee[1],
	    \qquad
	    [x,\lambda[1]]=(x\cdot\lambda)[1],\quad
	    [\lambda[1],\mu[1]]=0 .
	  \]
  The dual action is the coadjoint action:
  \[
    (x\cdot\lambda)(y)=-\lambda([x,y]).
  \]
  Choose a basis \(e_i\) of \(\mathfrak l\), dual basis
  \(\epsilon^i\), and structure constants
  \([e_i,e_j]=C^k_{ij}e_k\).  The CE coordinates are \(c^i\), dual to
  \(e_i\), and \(u_i\), dual to \(\epsilon^i[1]\), with
  \[
    |c^i|=1,\qquad |u_i|=0 .
  \]
  This is the cotangent-coordinate grading: \(u_i\) is the
  CE-suspended coordinate dual to the odd cotangent generator, not the
  ordinary graded dual computed from the BV degree of \(\epsilon^i[1]\).
  The polyvector model is
  \[
    \mathrm{PV}\bigl(\mathbf S(\mathfrak l)\bigr)
    =
    \mathbf S(O_i)\otimes\Lambda(\theta^i),
    \qquad |O_i|=0,\quad |\theta^i|=1,
  \]
  equipped with the Lichnerowicz differential for the linear Poisson
  tensor
  \[
    \pi=\frac12 C^k_{ij}O_k\theta^i\theta^j .
  \]
  The completed version is obtained by replacing tensor, symmetric, and
  exterior powers by their continuous completed analogues.
\end{defn}

\begin{thm}[Finite cotangent CE/PV identity]\label{thm:finite-cotangent-ce-pv-identity}
  Let \(\mathfrak l\) be finite-dimensional.  In the notation of
  Definition~\ref{defn:cotangent-ce-pv-datum}, the assignment
  \[
    \Phi(c^i)=\theta^i,\qquad \Phi(u_i)=O_i
  \]
  extends uniquely to an isomorphism of dg graded-commutative algebras
  with their degree-\((-1)\) cotangent/Schouten brackets
  \[
    \Phi\colon
    C^\bullet_{\mathrm{CE}}\bigl(\mathfrak l\ltimes\mathfrak l^\vee[1]\bigr)
    \xrightarrow{\ \sim\ }
    \mathrm{PV}\bigl(\mathbf S(\mathfrak l)\bigr).
  \]
  The differentials are
  \[
    d_{\mathrm{CE}}c^k=-\frac12 C^k_{ij}c^ic^j,\qquad
    d_{\mathrm{CE}}u_k=C^\ell_{kj}u_\ell c^j,
  \]
  and
  \[
    d_\pi\theta^k=-\frac12 C^k_{ij}\theta^i\theta^j,\qquad
    d_\pi O_k=C^\ell_{kj}O_\ell\theta^j.
  \]
\end{thm}

\begin{proof}
  The underlying graded algebras are freely generated by the same even
  and odd coordinate sets after the displayed substitution.  Hence
  \(\Phi\) is an isomorphism of graded-commutative algebras.

  The CE formula for \(c^k\) is the ordinary Chevalley--Eilenberg
  differential dual to \([e_i,e_j]=C^k_{ij}e_k\).  The formula for \(u_k\)
  is the dual of the coadjoint action in the semidirect product:
  \([e_k,e_j]=C^\ell_{kj}e_\ell\) makes \(e_j\) act on the coordinate
  \(u_k\) by \(C^\ell_{kj}u_\ell\).  On the PV side,
  \(d_\pi=[\pi,-]_S\) and the Schouten identities
  \[
    [\theta^i,O_k]_S=\delta^i_k,\qquad
    [\theta^i,\theta^j]_S=[O_i,O_j]_S=0
  \]
  give exactly the two displayed Lichnerowicz formulas.  Thus
  \(\Phi d_{\mathrm{CE}}=d_\pi\Phi\) on generators.  Both differentials
  are degree-one derivations, so equality propagates to all words.

  The degree-\((-1)\) cotangent/Schouten brackets are also checked on
  generators:
  \[
    \{c^i,u_j\}_{\mathrm{CE}}=\delta^i_j,\qquad
    [\theta^i,O_j]_S=\delta^i_j,
  \]
  with all \(c,c\), \(u,u\), \(\theta,\theta\), and \(O,O\) brackets
  equal to zero.  The biderivation rule gives the full bracket identity.
  The square-zero check on the cotangent coordinate is the coadjoint
  representation identity in coordinates:
  \[
    d_{\mathrm{CE}}^2u_k
    =
    C^\ell_{kj}C^m_{\ell r}u_m c^r c^j
    -\frac12 C^\ell_{kj}C^j_{rs}u_\ell c^r c^s
    =0,
  \]
  because
  \([\operatorname{ad}^{\vee}_{e_r},\operatorname{ad}^{\vee}_{e_s}]
    =\operatorname{ad}^{\vee}_{[e_r,e_s]}\).
\end{proof}

\begin{defn}[Formal-disk coordinate CE/PV completion]\label{defn:formal-disk-ce-pv-completion}
  Let \(H_d=\operatorname{Span}\{z_1^az_2^b:a+b=d\}\) and
  \[
    \mathfrak h=\prod_{d\geq1}H_d=(z_1,z_2)\C[[z_1,z_2]] .
  \]
  Its continuous dual is
  \(\mathfrak h^\vee_{\mathrm{cont}}=\bigoplus_{d\geq1}H_d^\vee\).
  Write \(c^{a,b}\) for the CE coordinate dual to
  \(H_{a,b}=z_1^az_2^b\), and \(u_{a,b}\) for the CE coordinate dual to
  the shifted cotangent generator \((H_{a,b})^\vee[1]\).  When
  \(I=(a,b)\) with \(a+b>0\), write
  \[
    H_I=z_1^az_2^b,\qquad
    \eta^I=(H_I)^\vee[1]\in\mathfrak h^\vee_{\mathrm{cont}}[1].
  \]

  For \(M,L\geq0\), let \(\mathcal B_{M,L}^{\mathrm{CE}}\) be the
  finite set of graded-commutative monomials in the letters
  \(u_{a,b},c^{a,b}\) of word length at most \(L\), all of whose labels
  satisfy \(1\leq a+b\leq M\).  The coordinate product algebra used here
  is the product over all finite words, equivalently the inverse limit of
  finite coordinate projections
  \[
    C^\bullet_{\mathrm{CE},\mathrm{coord}}(\mathfrak g)
    :=
    \widehat{\mathbf S}_{\mathrm{coord}}
    \bigl(\mathfrak g^\vee_{\mathrm{cont}}[-1]\bigr)
    =
    \varprojlim_{M,L}\prod_{w\in\mathcal B_{M,L}^{\mathrm{CE}}}\C w ,
    \qquad
    \mathfrak g=\mathfrak h\ltimes
      \mathfrak h^\vee_{\mathrm{cont}}[1].
  \]
	  The PV completion is defined by the same rule with \(u_{a,b}\)
	  replaced by \(O_{a,b}\) and \(c^{a,b}\) by \(\theta^{a,b}\):
	  \[
	    A_{\partial,\mathrm{coord}}
	    :=\widehat{\mathbf S}_{\mathrm{coord}}(\mathfrak h),
	    \qquad
	    \mathrm{PV}_{\mathrm{coord}}(\mathfrak h)
	    =
	    \mathrm{PV}_{\mathrm{coord}}(A_{\partial,\mathrm{coord}})
	    :=
	    \widehat{\mathbf S}_{\mathrm{coord}}(O_{a,b})
	    \widehat\otimes
    \widehat\Lambda_{\mathrm{coord}}(\theta^{a,b}).
  \]
	  These product algebras are coordinate enlargements.  They strictly
  contain the literal continuous CE cochains: for example an infinite
  length-one sum in the \(c^{a,b}\)-coordinates is a vector in the
  coordinate product but not a continuous linear functional on
  \(\mathfrak h=\prod_d H_d\).  A genuine completed \(P_0\)-object is the
  mixed continuous subspace in which the \(c/\theta\)-linear directions
  have continuous-dual support and the contraction pairing is a
  continuous biderivation.  Without this mixed condition the construction
	  is only an ambient pro-coordinate product.  The projections are finite
	  coordinate tests for the displayed differential formulas, not
	  Lie-algebra quotients of the full formal Hamiltonian algebra.
  A map out of either coordinate product is \emph{coordinate-continuous}
  when every output coefficient depends on finitely many input
  coefficients after projection to any finite set of word labels.
  Multiplication is defined by the usual finite decompositions of a fixed
  output word into two subwords; hence each output coefficient of a
  product is a finite sum.  The infinite behavior enters through the
  differential and contraction formulas, where coordinate-continuity or
  bracket-admissibility is the required summability condition.
		\end{defn}

\begin{defn}[Bracket- and kernel-admissible CE/PV realizations]\label{defn:bracket-kernel-admissible}
  Let \(A_{\partial,\mathrm{coord}}\) and
  \(\mathrm{PV}_{\mathrm{coord}}(A_{\partial,\mathrm{coord}})\) be as in
  Definition~\ref{defn:formal-disk-ce-pv-completion}.  A
  \emph{bracket-admissible realization} is a complete Hausdorff
  topological subalgebra
  \(i_B\colon B\hookrightarrow
  \mathrm{PV}_{\mathrm{coord}}(A_{\partial,\mathrm{coord}})\) with
  continuous inclusion such that multiplication
  \(B\widehat\otimes B\to B\), the restricted differential
  \(d_\pi^{\mathrm{coord}}\colon B\to B\), and the Schouten contraction
  bracket \(B\widehat\otimes B\to B[-1]\) are continuous.  The source
  \(\Phi_{\mathrm{coord}}^{-1}(B)\) is given the pullback topology; it is
  required to be closed under \(d_{\mathrm{CE}}\), multiplication, and the
  transported cotangent bracket, with these operations continuous.  For
  such \(B\), write
  \[
    \Phi^B_{P_0}\colon
    \Phi_{\mathrm{coord}}^{-1}(B)\xrightarrow{\sim}B
  \]
  for the restricted completed dg \(P_0\)-map.

  If \(S\subset\mathfrak h\) is a class of Hamiltonians, \(B\) is
  \emph{\(S\)-Hamiltonian-spanning} if \(O_f\in B\) and
  \(d_\pi^{\mathrm{coord}}O_f\in B\) for every \(f\in S\).  Boundary
  Hamiltonian statements inside \(\Phi^B_{P_0}\) are made only for such
  \(f\).

  A \emph{kernel-admissible datum} is a bracket-admissible \(B\) together
  with a specified completed tensor product, hence a completed
  convolution dg Lie algebra \(\mathcal K_B\), in which the diagonal sums
  \[
    \sum_I H_I\otimes\theta^I,\qquad
    \sum_I \eta^I\otimes O_I
  \]
  converge and the differential and convolution bracket are continuous.
  In that case set, as an element of \(\mathcal K_B\),
  \[
    \Theta_B=\sum_I H_I\otimes\theta^I+\sum_I\eta^I\otimes O_I .
  \]
  The strict ambient product/direct-sum completion is not
  kernel-admissible unless this diagonal convergence is supplied as
  additional structure.

	  The finite-support subalgebra \(B_{\mathrm{fin}}\) is a finite-test
	  model in every finite coordinate window; it is not a global
	  bracket-admissible \(B\) for the full formal disk unless the chosen
	  completion is also stable under the infinite coadjoint differential.
	  For \(q>1\) the weighted rapidly decreasing test target \(E_q\) is defined by the
	  seminorms
  \[
    p_n(x)=
    \sup_{a,b}\,|x_{a,b}|\,q^{a+b}(1+a+b)^n,\qquad n\geq0,
  \]
  and by the analogous finite-word seminorms on polyvectors.  The
	  monomial estimate
  \[
    |ad-bc|\leq (a+b)(c+d)
  \]
	  is only the first local convolution bound needed to absorb the
	  Poisson structure constants into the polynomial factor of the
	  seminorms.  Thus a Köthe completion built from these seminorms is a
	  candidate bracket-admissible \(B\) only after the chosen finite-word
	  tensor convention proves multiplication, the coordinate differential,
	  and the Schouten contraction continuous; it is kernel-admissible only
	  if the two diagonal sums above converge in that completed tensor
	  product.
\end{defn}

\begin{lem}[Formal-disk coordinate differential formulas]\label{lem:formal-disk-ce-schouten-continuity}
  In the coordinate product completion of
  Definition~\ref{defn:formal-disk-ce-pv-completion}, the formulas
  \[
    [H_{a,b},H_{c,d}]=(ad-bc)H_{a+c-1,b+d-1}
  \]
  with \(H_{r,s}=0\) for \(r<0\), \(s<0\), or \(r+s=0\), determine
  degree-one coordinate-continuous derivations
  \(d_{\mathrm{CE}}\) and \(d_\pi^{\mathrm{coord}}\), and these derivations
  satisfy \(d_{\mathrm{CE}}^2=(d_\pi^{\mathrm{coord}})^2=0\).  On generators,
  \[
    d_{\mathrm{CE}}c^{r,s}
      =
      -\frac12
      \sum_{\substack{a+c-1=r\\ b+d-1=s}}
      (ad-bc)c^{a,b}c^{c,d},
  \]
  \[
    d_{\mathrm{CE}}u_{r,s}
      =
      \sum_{c,d}
      (rd-sc)\,
      u_{r+c-1,s+d-1}c^{c,d},
  \]
    and the coordinate PV differential \(d_\pi^{\mathrm{coord}}\) has the identical formulas after
  \(c^{a,b}\mapsto\theta^{a,b}\) and \(u_{a,b}\mapsto O_{a,b}\).
  Terms whose output label is the reduced constant \(H_{0,0}\) are
  omitted.
\end{lem}

\begin{proof}
  For \(d_{\mathrm{CE}}c^{r,s}\), the equations
  \(a+c-1=r\) and \(b+d-1=s\) have finitely many non-negative
  solutions, hence the displayed sum is finite.

  The coadjoint term \(d_{\mathrm{CE}}u_{r,s}\) is infinite before
  projection.  Its projection to the finite coordinate window
  \(\mathcal B_{M,L}^{\mathrm{CE}}\) keeps only terms with
  \(c+d\leq M\) and \(r+s+c+d-2\leq M\).  There are finitely many such
  pairs \((c,d)\).  The word length increases by one, so the length
  projection is finite as well.  Thus every finite coordinate
  projection of \(d_{\mathrm{CE}}u_{r,s}\) is finite, which is exactly
  continuity in the product topology.

  A completed monomial has finite word length before passage to the
  product completion, so the derivation rule applies to finitely many
  factors.  Continuity extends the derivation uniquely to arbitrary
  products over the finite-word basis.  The PV coordinate formulas have
  the same coefficient support, with \(O,\theta\) in place of \(u,c\), so
  the same argument proves continuity of \(d_\pi^{\mathrm{coord}}\).

  The square-zero assertions are checked before completion and then
  extended by continuity.  The CE formulas are the continuous dual of the
  genuine Hamiltonian Lie bracket on
  \(\C[[z_1,z_2]]/\C\cdot1\), equivalently on the zero-constant
  representatives \((z_1,z_2)\C[[z_1,z_2]]\), whose Jacobi identity gives
  \(d_{\mathrm{CE}}^2c^{r,s}=0\).  On the cotangent coordinates,
  \[
    d_{\mathrm{CE}}^2u_{r,s}
    =
    C^L_{(r,s)J}C^M_{LR}u_Mc^Rc^J
    -\frac12 C^L_{(r,s)J}C^J_{RS}u_Lc^Rc^S
    =0,
  \]
  where the constants \(C^K_{IJ}\) are those of
  \([H_I,H_J]=C^K_{IJ}H_K\), and the equality is exactly
  \([\operatorname{ad}^{\vee}_{H_R},\operatorname{ad}^{\vee}_{H_S}]
    =\operatorname{ad}^{\vee}_{[H_R,H_S]}\).
  On the PV side the displayed coordinate formulas define a derivation
  \(d_\pi^{\mathrm{coord}}\); its square on generators is the same Jacobi
  and coadjoint identity with \(O,\theta\) replacing \(u,c\).  This avoids
  asserting that the infinite
  linear Poisson bivector is a globally summable Schouten element in the
  strict ambient product.  Since both differentials are continuous
  coordinate derivations, generator-level square-zero implies
  square-zero on the completed coordinate product algebras.
\end{proof}

\begin{thm}[Coordinate formal-disk CE/PV theorem and admissible \(P_0\) realizations]\label{thm:formal-disk-completed-ce-pv}
  Let
  \[
    \mathfrak h=\C[[z_1,z_2]]/\C\cdot1
      \cong (z_1,z_2)\C[[z_1,z_2]]
  \]
  where the isomorphism chooses the zero-constant representative, with
  the projected Poisson bracket modulo constants
  \[
    \{f,g\}=\partial_{z_1}f\,\partial_{z_2}g
      -\partial_{z_2}f\,\partial_{z_1}g .
  \]
  Complete the coordinate CE and PV algebras as in
  Definition~\ref{defn:formal-disk-ce-pv-completion}.  Then
  \[
	    \Phi_{\mathrm{coord}}\colon
	    C^\bullet_{\mathrm{CE},\mathrm{coord}}
	    \bigl(\mathfrak h\ltimes\mathfrak h^\vee_{\mathrm{cont}}[1]\bigr)
    \xrightarrow{\ \sim\ }
    \mathrm{PV}_{\mathrm{coord}}(\mathfrak h),
    \qquad
    c^{a,b}\mapsto\theta^{a,b},\quad u_{a,b}\mapsto O_{a,b},
	  \]
	  is an isomorphism of completed coordinate dg commutative algebras.
	  If \(B\subset
	  \mathrm{PV}_{\mathrm{coord}}(A_{\partial,\mathrm{coord}})\) is
	  bracket-admissible in the sense of
	  Definition~\ref{defn:bracket-kernel-admissible}, then
	  \[
	    \Phi^B_{P_0}\colon\Phi_{\mathrm{coord}}^{-1}(B)
	    \xrightarrow{\sim}B
	  \]
	  is an isomorphism of completed dg \(P_0\)-algebras.  If \(B\) is
	  kernel-admissible as part of a datum \((B,\mathcal K_B)\), then the
	  diagonal element
	  \[
	    \Theta_B=\sum_I H_I\otimes\theta^I+\sum_I\eta^I\otimes O_I
	  \]
	  satisfies the Maurer--Cartan equation
	  \(d\Theta_B+\frac12[\Theta_B,\Theta_B]=0\) in \(\mathcal K_B\).  The
	  full product
	  coordinate algebra is not asserted to carry the global continuous
	  contraction bracket or a strict kernel; the missing infinite diagonal
	  is the strict Casimir obstruction of
	  Theorem~\ref{thm:lane-strict-casimir-obstruction}.  The admissible
	  formal-local interface is
	  \[
	    \mathfrak U_{\mathfrak h}
	    =(\Phi_{\mathrm{coord}},\{\Phi^B_{P_0}\}_B,\{\Theta_B\}_B),
	  \]
	  with \(B\) ranging over bracket-admissible realizations in the middle
	  term and kernel-admissible data in the last term.  At the coordinate
	  dg-algebra level, and inside \(\Phi^B_{P_0}\) for every
	  Hamiltonian-spanning \(B\), the stable boundary evaluation of
	  Theorem~\ref{thm:bulk-boundary} sends the degree-zero cotangent
	  coordinate \(u_f\) to the boundary Hamiltonian
  \[
    B_f=\overline{\operatorname{Tr}f(\phi_1,\phi_2)} ,
  \]
  while \(c^f\) maps to the Hamiltonian vector field of \(B_f\).  This
  is the cochain-level bulk-boundary Koszul shadow; it does not use
  lower-central Tate-pronilpotence of the full Hamiltonian Lie algebra.
\end{thm}

\begin{proof}
	  The completed coordinate source and target are products over the same
	  finite-word coordinate set after the substitutions
		  \(c^{a,b}\leftrightarrow\theta^{a,b}\) and
		  \(u_{a,b}\leftrightarrow O_{a,b}\).  Hence
	  \(\Phi_{\mathrm{coord}}\) is a continuous
  isomorphism of completed graded-commutative algebras, with continuous
  inverse given by the reverse substitution.

  Lemma~\ref{lem:formal-disk-ce-schouten-continuity} gives
  coordinate-continuous
  derivations \(d_{\mathrm{CE}}\) and \(d_\pi^{\mathrm{coord}}\).  On the generators
  \(c^{r,s}\) and \(u_{r,s}\), the displayed CE formulas are the duals
  of the Hamiltonian bracket and coadjoint action in
  \(\mathfrak h\ltimes\mathfrak h^\vee_{\mathrm{cont}}[1]\).  The
  PV coordinate formulas are the coefficientwise Lichnerowicz formulas
  for the linear Poisson tensor
  \[
    \pi=\frac12 C^K_{IJ}O_K\theta^I\theta^J .
  \]
	  Therefore \(\Phi_{\mathrm{coord}} d_{\mathrm{CE}}=d_\pi^{\mathrm{coord}}\Phi_{\mathrm{coord}}\) on \(c^{r,s}\) and
	  \(u_{r,s}\).  Both sides are continuous degree-one derivations, so the
  equality extends from generators to every finite word and then, by the
  product completion, to the whole completed algebra.

	  On the finite-support mixed subalgebra the cotangent/Schouten bracket is
  checked on generators:
  \[
    \{c^{a,b},u_{r,s}\}_{\mathrm{CE}}
      =\delta_{a,r}\delta_{b,s},
    \qquad
    [\theta^{a,b},O_{r,s}]_S
      =\delta_{a,r}\delta_{b,s},
  \]
  and all \(c,c\), \(u,u\), \(\theta,\theta\), and \(O,O\) brackets
  vanish.  This gives bracket compatibility on finite words and on any
  mixed topology in which the \(c/\theta\)-directions have finite label
  support.  It does not extend to the arbitrary product completion: the
  bracket of \(\sum_I c^I\) with \(\sum_I u_I\), which are both admitted
  by the coordinate product, would require the divergent diagonal
  coefficient \(\sum_I1\).  Thus the full product statement above is a
	  dg coordinate reconstruction, not a global completed \(P_0\)-algebra
	  theorem.

	  If \(B\) is bracket-admissible, the preceding generator calculation
	  restricts to \(B\), and all infinite contractions that occur in the
	  \(P_0\) bracket converge by definition of \(B\).  Hence the restricted
	  map \(\Phi^B_{P_0}\) is a completed dg \(P_0\)-isomorphism.  If \(B\)
	  is kernel-admissible, the two diagonal sums defining \(\Theta_B\)
	  converge in the specified convolution dg Lie algebra \(\mathcal K_B\).
	  The Maurer--Cartan equation is
	  then checked coefficientwise in the completed convolution complex:
	  every requested coefficient sees only the finitely many labels
	  contributing to that coefficient, and the equality is exactly the
	  Jacobi and coadjoint identity in the full Hamiltonian Lie algebra.
	  Convergence in \(B\) passes these coefficient identities to the
	  completed tensor product.  No finite coordinate window is being used
	  as a Hamiltonian Lie quotient.  Without diagonal
	  convergence this sentence has no target object, which is precisely the
	  strict Casimir obstruction.

		  The coordinate moment-map shadow is the restriction of
	  \(\Phi_{\mathrm{coord}}\) to the cotangent
  coordinates: \(u_f\mapsto O_f\).  The Kirillov--Poisson bracket gives
  \(\{O_f,O_g\}=O_{\{f,g\}}\).  Theorem~\ref{thm:bulk-boundary} sends
  \(O_f\) to
  \(\overline{\operatorname{Tr}f(\phi_1,\phi_2)}\) and preserves the
  Hamiltonian bracket in the stable connected sector.  Inside a
  bracket-admissible \(P_0\)-realization this sentence applies only for
  the \(f\)'s in the Hamiltonian-spanning set of \(B\).  Thus \(u_f\) is
  the bulk source of the boundary Hamiltonian, while \(c^f\mapsto\theta^f\)
  records the corresponding Hamiltonian vector field.
\end{proof}

\begin{thm}[Abstract bracket-compatible CE/PV recognition criterion]\label{thm:universal-formal-local-ce-pv-recognition}
  Let \(L=\varprojlim_\alpha L_\alpha\) be a complete Hausdorff
  filtered dg Lie algebra over \(\C\), with finite-dimensional
  bracket-compatible dg Lie quotients \(L_\alpha\), and let
  \(D=L^\vee_{\mathrm{cont}}=\varinjlim_\alpha L_\alpha^\vee\).
  Put \(G=L\ltimes D[1]\).  Form the completed coordinate CE and PV
  algebras by taking products over finite words and finite coordinate
  windows.  Assume:
  \begin{enumerate}[label=(U\arabic*)]
    \item the internal differential, Lie bracket on \(L\), and
      coadjoint action on \(D\) are continuous, equivalently every
      finite coordinate projection of \(d_Lx\), \([x,y]\), and
      \(x\cdot\lambda\) factors through a finite Lie quotient of
      the inputs;
    \item the linear Poisson bivector
      \[
        \pi_L=\frac12 C^k_{ij}O_k\theta^i\theta^j
      \]
      determined by the structure constants of \(L\) is a continuous
      Schouten bivector on
      \(\widehat{\mathbf S}_{\mathrm{coord}}(L)\);
    \item the source cotangent bracket
      \(\{c^i,u_j\}_{\mathrm{CE}}=\delta^i_j\) and the target Schouten
      contraction bracket \([\theta^i,O_j]_S=\delta^i_j\) are continuous
      biderivations on the chosen mixed or weighted completion.  This is
      a summability hypothesis, not a consequence of the generator
      formula on the ambient product.
  \end{enumerate}
  Then the generator assignment
  \[
    \Phi_L\colon
    C^\bullet_{\mathrm{CE},\mathrm{cont}}(G)
    \xrightarrow{\ \sim\ }
    \mathrm{PV}\bigl(\widehat{\mathbf S}_{\mathrm{coord}}(L)\bigr),
    \qquad
    c^i\mapsto\theta^i,\quad u_i\mapsto O_i,
  \]
  is an isomorphism of completed dg algebras with the degree-\((-1)\)
  cotangent/Schouten contraction bracket.  Without the summability
  condition in (U3), the conclusion is only the compatible pro
  finite-window dg-commutative CE/PV identity and the global contraction
  bracket is obstruction data.  This is a formal-local cochain theorem.
  It does not require lower-central Tate-pronilpotence and it does not assert
  that the coordinate pairing is represented by a strict kernel
  \(C_L\in L\widehat\otimes D\).
  The strict formal Hamiltonian disk is not an instance of this theorem
  unless such a bracket-compatible finite quotient system is supplied;
  for the disk itself one uses the coordinate finite-test theorem above.
\end{thm}

\begin{proof}
  The completed CE algebra of \(G=L\ltimes D[1]\) and the completed PV
  algebra of the linear formal space \(L\) are products over the same
  finite-word coordinate monomials after \(c^i\leftrightarrow\theta^i\)
  and \(u_i\leftrightarrow O_i\).  Thus \(\Phi_L\) is a continuous
  graded-commutative algebra isomorphism with continuous inverse.

  Assumption (U1) makes the CE differential continuous: every finite
  output coordinate of the internal differential, of
  \(d_{\mathrm{CE}}c^k\), and of \(d_{\mathrm{CE}}u_i\) is computed from
  finitely many input coordinates.  The internal differential is carried
  to the identical linear vector field on the PV side by the coordinate
  substitution.  Assumption (U2) makes \(d_{\pi_L}=[\pi_L,-]_S\) a
  continuous differential on the PV side.  On generators the bracket
  terms are the same structure-constant formulas:
  \[
    d_{\mathrm{CE}}c^k=-\frac12 C^k_{ij}c^ic^j,\qquad
    d_{\pi_L}\theta^k=-\frac12 C^k_{ij}\theta^i\theta^j,
  \]
  and
  \[
    d_{\mathrm{CE}}u_i=C^\ell_{ij}u_\ell c^j,\qquad
    d_{\pi_L}O_i=C^\ell_{ij}O_\ell\theta^j .
  \]
  Jacobi for \(L\) gives \(d_{\mathrm{CE}}^2=0\) and
  \([\pi_L,\pi_L]_S=0\), hence \(d_{\pi_L}^2=0\).  Since both
  differentials are continuous derivations, generator-level
  compatibility gives \(\Phi_Ld_{\mathrm{CE}}=d_{\pi_L}\Phi_L\) on the
  completed algebras.

  Assumption (U3) gives the same result for the contraction brackets:
  \[
    \{c^i,u_j\}_{\mathrm{CE}}=\delta^i_j
    \quad\text{and}\quad
    [\theta^i,O_j]_S=\delta^i_j
  \]
  on generators, with all same-type brackets zero.  Continuity and the
  biderivation rule extend this equality to finite words and then to the
  chosen mixed or weighted completion.  The last sentence records the
  boundary of the theorem: a strict Casimir kernel is stronger than the
  coordinate biderivation, while the unweighted ambient product may fail
  even to carry that biderivation globally.
\end{proof}

\begin{prop}[Shifted-cotangent coordinate is the boundary Hamiltonian]\label{prop:cotangent-coordinate-boundary-hamiltonian}
  In the completed formal-disk identity of
  Theorem~\ref{thm:formal-disk-completed-ce-pv}, the boundary Hamiltonian
  \(O_f\) is the image of the shifted-cotangent CE coordinate \(u_f\):
  \[
	    \Phi_{\mathrm{coord}}(u_f)=O_f.
  \]
  The Hamiltonian CE coordinate \(c^f\) maps to the vector-field
  generator \(\theta^f\), not to \(O_f\).  Equivalently,
  \[
	    d_\pi^{\mathrm{coord}} O_f=X_{O_f}=\Phi_{\mathrm{coord}}(d_{\mathrm{CE}}u_f).
  \]
  After stable boundary evaluation, \(O_f\) becomes
  \(\overline{\operatorname{Tr}f(\phi_1,\phi_2)}\), the moment-map
  shadow of the Hamiltonian \(f\) on the matrix probe.
\end{prop}

\begin{proof}
  This is the \(u\)-generator part of the defining formula for
  \(\Phi_{\mathrm{coord}}\).
  For \(f=H_k\), the CE differential is
  \(d_{\mathrm{CE}}u_k=C^\ell_{kj}u_\ell c^j\).  Applying
  \(\Phi_{\mathrm{coord}}\) gives
  \(C^\ell_{kj}O_\ell\theta^j\), which is exactly the Hamiltonian vector
	  field \(d_\pi^{\mathrm{coord}}O_k\); on a bracket-admissible \(B\) this
	  is the Schouten bracket \([\pi,O_k]_S\).  Thus the cotangent coordinate, not the
  Hamiltonian ghost coordinate, carries the boundary Hamiltonian.  The
  last assertion is Theorem~\ref{thm:bulk-boundary} applied to the
  degree-zero Kirillov map \(f\mapsto O_f\).
\end{proof}

\begin{thm}[Strict product/direct-sum Casimir obstruction]\label{thm:lane-strict-casimir-obstruction}
  Let
  \[
    \mathfrak h_\Pi=\prod_{d\geq1}H_d,\qquad
    D_\oplus=\bigoplus_{d\geq1}H_d^\vee
  \]
  be the strict product/direct-sum coefficient pair of the formal
  Hamiltonian disk.  There is no tensor
  \(K\in\mathfrak h_\Pi\widehat\otimes D_\oplus\) whose projection to
  every finite window \(d\leq M\) is the finite Casimir
  \[
    C_{\leq M}=\sum_{1\leq d\leq M}\sum_i e_{d,i}\otimes e^i_{d}.
  \]
  The same statement holds with tensor factors reversed.
\end{thm}

\begin{proof}
  In the strict tensor convention for the product/direct-sum Tate pair,
  an element of \(\mathfrak h_\Pi\widehat\otimes D_\oplus\) has finite
  support in the direct-sum factor \(D_\oplus\).  Thus there is a degree
  \(M_0\) such that all \(H_d^\vee\)-components of \(K\) vanish for
  \(d>M_0\).  Projecting to the window \(d\leq M_0+1\) then gives zero in
  the \(H_{M_0+1}\otimes H_{M_0+1}^\vee\) component, while
  \(C_{\leq M_0+1}\) has the nonzero identity tensor in that component.
  This contradiction proves the claim.  Reversing the tensor factors is
  identical.
\end{proof}

\begin{thm}[Primitive-label obstruction to Matlis cotangent-current equivariance]\label{thm:lane-one-psi-not-cotangent-current}
  Let \(V_\Psi\) be the span of primitive PBW special-fibre labels
  \(\widetilde\Psi_{a,b}\), \(a+b>0\), with the reduced scalar target
  killed and with the linear Hamiltonian action
  \[
    O_{1,0}\cdot\widetilde\Psi_{a,b}
      =b\,\widetilde\Psi_{a,b-1},\qquad
    O_{0,1}\cdot\widetilde\Psi_{a,b}
      =-a\,\widetilde\Psi_{a-1,b}.
  \]
  Let \(\mathcal P\) be the principal-part module with
  \[
    z_1\cdot\rho_{a,b}=-(b+1)\rho_{a,b+1},\qquad
    z_2\cdot\rho_{a,b}=(a+1)\rho_{a+1,b}.
  \]
  There is no \(\mathfrak h\)-module map
  \(F\colon V_\Psi\to\mathcal P\) satisfying
  \(F(\widetilde\Psi_{a,b})=\lambda_{a,b}\rho_{a,b}\) with all
  \(\lambda_{a,b}\neq0\).
\end{thm}

\begin{proof}
  Apply equivariance for \(z_1\) to \(\widetilde\Psi_{a,0}\), \(a>0\).  On the
  PBW special-fibre side,
  \(O_{1,0}\cdot\widetilde\Psi_{a,0}=0\).  Hence
  \(F(O_{1,0}\cdot\widetilde\Psi_{a,0})=0\).  On the principal-part side,
  \[
    z_1\cdot F(\widetilde\Psi_{a,0})
    =
    z_1\cdot(\lambda_{a,0}\rho_{a,0})
    =
    -\lambda_{a,0}\rho_{a,1},
  \]
  which is nonzero.  This contradicts equivariance.  The \(z_2\)-action
  gives the analogous contradiction at \(a=0\).
\end{proof}

\begin{prop}[Low-degree obstruction to a pronilpotent full-disk splitting]\label{prop:lane-low-degree-pronilpotent-obstruction}
  The subspace \(\mathfrak m^3\subset\mathfrak h=(z_1,z_2)\C[[z_1,z_2]]\)
  is not an ideal of the full reduced Hamiltonian Lie algebra
  \(\mathfrak h\).  Consequently the vector-space decomposition
  \[
    \mathfrak h=(H_1\oplus H_2)\oplus\mathfrak m^3,
    \qquad
    H_d=\operatorname{Span}\{z_1^az_2^b:a+b=d\},
  \]
  does not make \(\mathfrak h\) a semidirect product of the low-degree
  sector with a pronilpotent ideal.
\end{prop}

\begin{proof}
  The linear Hamiltonian \(z_1\in H_1\) lowers \(z_2^3\in\mathfrak m^3\):
  \[
    \{z_1,z_2^3\}=3z_2^2\in H_2.
  \]
  Since \(H_2\cap\mathfrak m^3=0\), the bracket does not lie in
  \(\mathfrak m^3\).  Thus \(\mathfrak m^3\) is not an ideal in
  \(\mathfrak h\).  A semidirect-product presentation with
  pronilpotent radical \(\mathfrak m^3\) would require that radical to
  be an ideal, so no such presentation exists.
\end{proof}

\begin{thm}[Relative filtered CE/PV reconstruction criterion]\label{thm:relative-filtered-koszul-lift}
  Let \(\mathfrak h_{\mathrm{rel}}\) be a complete filtered dg Lie algebra
  with a finite-dimensional reductive degree-zero subalgebra
  \(\mathfrak a\), a complete pronilpotent dg ideal \(\mathfrak n\), and
  a filtered splitting
  \[
    \mathfrak h_{\mathrm{rel}}\cong\mathfrak a\ltimes\mathfrak n .
  \]
  Let \(D_{\mathrm{rel}}\) be a continuous \(\mathfrak h_{\mathrm{rel}}\)-module
  with a perfect filtered pairing with \(\mathfrak h_{\mathrm{rel}}\) on
  every finite quotient, and assume the coevaluation tensor belongs to
  \(\mathfrak h_{\mathrm{rel}}\widehat\otimes D_{\mathrm{rel}}\).  Put
  \[
    \mathfrak g_{\mathrm{rel}}
      =\mathfrak h_{\mathrm{rel}}\ltimes D_{\mathrm{rel}}[1],
    \qquad
    A_{\mathrm{rel}}=\widehat{\mathbf S}(\mathfrak h_{\mathrm{rel}}).
  \]
  Then the generator assignment \(c\mapsto\theta\), \(u\mapsto O\)
  defines a filtered isomorphism
  \[
    C^\bullet_{\mathrm{CE},\mathrm{cont}}(\mathfrak g_{\mathrm{rel}})
    \xrightarrow{\ \sim\ }
    \mathrm{PV}(A_{\mathrm{rel}})
  \]
  of completed dg algebras with the degree-\((-1)\)
  cotangent/Schouten contraction bracket.  If, in addition, the relative
  bar-cobar functor is taken in the category of
  \(\mathfrak a\)-modules and the \(\mathfrak n\)-coalgebra is
  conilpotent, this cochain identity is promoted to relative
  Lie-Koszul duality:
  \[
    C^\bullet_{\mathrm{CE},\mathrm{cont}}(\mathfrak g_{\mathrm{rel}},\mathfrak a)
    \simeq
    \bigl(\widehat U_{\mathfrak a}(\mathfrak g_{\mathrm{rel}})\bigr)^!.
  \]
\end{thm}

\begin{proof}
  The cochain-level statement is the completed version of
  Theorem~\ref{thm:finite-cotangent-ce-pv-identity}: finite quotients of
  the filtration carry the finite CE/PV identity, the transition maps
  preserve \(c\mapsto\theta\), \(u\mapsto O\), and completeness identifies
  the inverse limit with the displayed completed complexes.  The
  coevaluation hypothesis is the exact condition needed to keep the
  cotangent generators inside the same completed tensor category.

  For the Koszul-dual statement, work relative to the reductive
  subalgebra \(\mathfrak a\).  The coalgebra generated by the
  pronilpotent ideal \(\mathfrak n\) is conilpotent, so the
  Lie/commutative bar-cobar equivalence applies to the relative CE
  chains.  Reductivity of \(\mathfrak a\) makes passage to
  \(\mathfrak a\)-modules exact on finite quotients; in the filtered
  inverse limit this gives the completed relative equivalence.  Dualizing
  yields the displayed cochain Koszul-dual form.

  This is a criterion, not a statement about the unmodified formal disk:
  Proposition~\ref{prop:lane-low-degree-pronilpotent-obstruction} shows
  that the actual decomposition
  \(\mathfrak h=(H_1\oplus H_2)\oplus\mathfrak m^3\) has no
  pronilpotent ideal \(\mathfrak m^3\) for the full bracket, so the
  criterion applies only after an independent relative model supplies the
  displayed semidirect-product hypotheses.
\end{proof}

\begin{lem}[Filtered cobar acceptance criterion]\label{lem:filtered-cobar-qiso-criterion}
  Let \(C\) be a coaugmented conilpotent dg coalgebra and \(B\) a
  complete augmented filtered \(A_\infty\)-algebra in the same complete
  linear category.  For a continuous degree-one map
  \(\tau\colon C\to B\), define
  \[
    \operatorname{MC}(\tau)
      =
      d_B\tau+\tau d_C
      +\sum_{r\geq2}m_r^B(\tau^{\otimes r})\Delta_C^{(r)}
      \in\Hom^2_{\mathrm{cont}}(C,B),
  \]
  where \(\Delta_C^{(r)}\) is the \(r\)-fold reduced coproduct.  The
  equation \(\operatorname{MC}(\tau)=0\) is equivalent to the existence
  of the adjoint complete cobar morphism
  \[
    \kappa_\tau\colon\Omega C\longrightarrow B .
  \]
  Assume the filtrations on \(\Omega C\) and \(B\) are bounded below,
  complete, separated, and strongly convergent, with finite-dimensional
  associated graded pieces in every total degree.  Then
  \(\kappa_\tau\) is an admissible filtered quasi-isomorphism if and
  only if
  \[
    H^\bullet\!\left(\operatorname{Cone}(
      \operatorname{gr}\kappa_\tau)\right)=0
  \]
  and the two finite-window Milnor defects
  \[
    \varprojlim\nolimits^1_M
      H^{n-1}\!\left((\Omega C)_{\leq M}\right),
    \qquad
    \varprojlim\nolimits^1_M
      H^{n-1}\!\left(B_{\leq M}\right)
  \]
  vanish for every \(n\).
\end{lem}

\begin{proof}
  The Maurer--Cartan formula is the bar-cobar adjunction written in the
  continuous convolution complex.  For a dg target it reduces to
  \(d_B\tau+\tau d_C+\mu_B(\tau\otimes\tau)\Delta_C=0\), which says
  exactly that \(\Omega C\to B\) commutes with differentials and
  multiplication.  The displayed \(A_\infty\) formula is the same
  statement with all target multiplications retained.

  Filter source and target by tensor length and by finite coordinate
  windows.  Exactness of the associated-graded cone is equivalent to
  \(\operatorname{gr}\kappa_\tau\) being a quasi-isomorphism.  The two
  \(\varprojlim^1\)-terms are the Milnor errors in computing the
  cohomology of the completed source and target from their finite
  windows.  With those terms zero, strong convergence gives the
  spectral-sequence comparison theorem in both directions.  Necessity is
  the definition of admissible filtered quasi-isomorphism in this
  envelope: it is detected on associated graded and on the exact
  finite-window towers.
\end{proof}

\begin{defn}[Universal open-trace \(A_\infty\)-Koszul acceptance datum]\label{defn:stable-ainfty-koszul-acceptance-datum}
  A universal open-trace \(A_\infty\)-Koszul acceptance datum consists of:
  \begin{enumerate}[label=(K\arabic*)]
    \item a completed augmented filtered \(A_\infty\)-algebra
      \(A_{\mathrm{op}}\) in the scalar-reduced stable trace/Koszul/BV
      sector, with associated graded
      \[
        \operatorname{gr}A_{\mathrm{op}}\cong
        \widehat{\mathbf S}(V_H),
        \qquad
        \operatorname{Indec}\operatorname{gr}A_{\mathrm{op}}
          =\mathfrak m/\mathfrak m^2\cong V_H,
      \]
      where the connected/primitive expression
      \[
        \pi_{\mathrm{Ham}}
        \operatorname{Conn}_{N\to\infty}
        \operatorname{Prim}_{\mathrm{st}}
        \bigl(\mc R_N^{GL_N},Q\bigr)
      \]
      is only this linear associated-graded symbol, not the algebra object
      itself;
    \item the completed closed target
      \[
        B_{\mathrm{CE/PV}}
          =
        C^\bullet_{\mathrm{CE},\mathrm{cont}}(\mathfrak g_{\mathrm{adm}})
          \cong
        \mathrm{PV}\bigl(
          \widehat{\mathbf S}_{\mathrm{coord}}(\mathfrak h_{\mathrm{adm}})
        \bigr),
      \]
      with the generator dictionary \(c^f\mapsto\theta^f\),
      \(u_f\mapsto O_f\);
    \item the open bar coalgebra
      \(C_{\mathrm{op}}=\mathrm{Bar}(A_{\mathrm{op}})\), equipped with a
      complete bounded-below filtration compatible with the bar
      differential, coproduct, and finite coordinate windows;
    \item a continuous degree-one map
      \[
        \tau_{\mathrm{op}}\colon
        C_{\mathrm{op}}\longrightarrow B_{\mathrm{CE/PV}}
      \]
      satisfying the explicit Maurer--Cartan equation
      \[
        d_{B_{\mathrm{CE/PV}}}\tau_{\mathrm{op}}
        +\tau_{\mathrm{op}} b_{\mathrm{op}}
        +\mu_{B_{\mathrm{CE/PV}}}
          (\tau_{\mathrm{op}}\otimes\tau_{\mathrm{op}})\Delta
        =0;
      \]
    \item the induced completed cobar morphism
      \[
        \kappa_{\tau_{\mathrm{op}}}\colon
        \Omega C_{\mathrm{op}}
        \longrightarrow
        B_{\mathrm{CE/PV}}
      \]
      whose associated-graded linear symbol is
      \[
        B_f\longmapsto O_f,\qquad
        \Theta_{\rho_{a,b}}\longmapsto\theta^{a,b};
      \]
    \item the acceptance obstruction complex
      \[
        \begin{aligned}
        \mathfrak C_{\tau_{\mathrm{op}}}
          &=
        \operatorname{Cone}(
          \operatorname{gr}\kappa_{\tau_{\mathrm{op}}})\\
          &\quad{}\oplus
        \bigoplus_n\varprojlim\nolimits^1_M
          H^{n-1}\!\left((\Omega C_{\mathrm{op}})_{\leq M}\right)[-n]\\
          &\quad{}\oplus
        \bigoplus_n\varprojlim\nolimits^1_M
          H^{n-1}\!\left((B_{\mathrm{CE/PV}})_{\leq M}\right)[-n].
        \end{aligned}
      \]
    \item exact continuous dualization on the same finite-window tower.
  \end{enumerate}
\end{defn}

\begin{thm}[Universal open-side filtered \(A_\infty\)-Koszul acceptance]\label{thm:stable-ainfty-koszul-under-hypotheses}
  For a datum in the sense of
  Definition~\ref{defn:stable-ainfty-koszul-acceptance-datum}, the following
  are equivalent:
  \begin{enumerate}[label=(\alph*)]
    \item \(\tau_{\mathrm{op}}\) is an accepted open-side twisting
      cochain, i.e. \(\kappa_{\tau_{\mathrm{op}}}\) is an admissible
      filtered quasi-isomorphism of complete augmented
      \(A_\infty\)-algebras with the displayed associated-graded symbol;
    \item the Maurer--Cartan equation in (K4) holds and the obstruction
      complex \(\mathfrak C_{\tau_{\mathrm{op}}}\) is exact.
  \end{enumerate}
  Under these equivalent conditions,
  \[
    \Omega\mathrm{Bar}(A_{\mathrm{op}})
      \xrightarrow{\ \sim\ }
    B_{\mathrm{CE/PV}}
  \]
  and hence \(A_{\mathrm{op}}\simeq B_{\mathrm{CE/PV}}\) in the
  completed filtered \(A_\infty\) homotopy category.  If the CE-side
  cobar map
  \[
    \Omega C_*^{\mathrm{CE}}(\mathfrak g_{\mathrm{adm}})
      \longrightarrow A_{\mathrm{op}}
  \]
  satisfies the same exact cone and \(\varprojlim^1\) test, exact
  continuous dualization gives the Koszul-dual cochain form
  \[
    C^\bullet_{\mathrm{CE},\mathrm{cont}}(\mathfrak g_{\mathrm{adm}})
      \simeq
    A_{\mathrm{op}}^! .
  \]
\end{thm}

\begin{proof}
  Apply Lemma~\ref{lem:filtered-cobar-qiso-criterion} with
  \(C=C_{\mathrm{op}}\) and \(B=B_{\mathrm{CE/PV}}\).  The
  Maurer--Cartan equation in (K4) is precisely the condition that
  \(\kappa_{\tau_{\mathrm{op}}}\) is a morphism.  Exactness of
  \(\mathfrak C_{\tau_{\mathrm{op}}}\) says simultaneously that
  \(\operatorname{gr}\kappa_{\tau_{\mathrm{op}}}\) is a
  quasi-isomorphism and that both finite-window towers compute completed
  cohomology.  The filtered comparison theorem then gives the accepted
  quasi-isomorphism, and the converse is the admissible-filtered
  definition.

  The bar-cobar counit
  \(\Omega\mathrm{Bar}(A_{\mathrm{op}})\to A_{\mathrm{op}}\) is a
  filtered quasi-isomorphism under the same boundedness and convergence
  hypotheses, so the accepted \(\kappa_{\tau_{\mathrm{op}}}\) identifies
  \(A_{\mathrm{op}}\) with \(B_{\mathrm{CE/PV}}\) in the homotopy
  category.  The final displayed equivalence is the CE-side specialization
  of the same criterion, followed by exact continuous dualization.
\end{proof}

\begin{stmt}[Statement: Matlis principal-part cotangent module]\label{thm:lane-principal-parts}
  \emph{Status.} \emph{proved degreewise stable}.  The proof uses
  residue calculus and Matlis local-cohomology duality for the
  \(m\)-adic topology on \(R=\C[[z_1,z_2]]\), after quotienting
  constants.

  \emph{Exact hypotheses.} Work over \(R=\C[[z_1,z_2]]\) with maximal
  ideal \(\mathfrak m=(z_1,z_2)\); put
  \(\mathfrak h=R/\C\cdot1\); take continuous duals for the
  \(m\)-adic topology; give \(\mathcal P\) the locally finite
  direct-sum topology transported from
  \(\mathfrak h^\vee_{\mathrm{cont}}\); normalize the residue pairing
  only up to one nonzero scalar; and, for the Fourier--Rees comparison
  only, distinguish the Fourier delta Rees lattice
  \(\mathcal D_\hbar/(z_1,z_2)\) from the full \v{C}ech
  local-cohomology lattice before \(\hbar\) is inverted.

  \emph{Local assertion.} The continuous dual of the reduced Hamiltonian
  Lie algebra is the reduced Matlis-local-cohomology module
  \[
    \mathfrak h^\vee_{\mathrm{cont}}
    \cong
    \mathcal P
    =
    \operatorname{Ann}_{\mathrm{res}}(\C\cdot1)
    \subset
    H^2_{\mathfrak m}(\C[[z_1,z_2]])\,dz_1dz_2,
  \]
  with basis
  \(\rho_{a,b}=z_1^{-a-1}z_2^{-b-1}\,dz_1dz_2\), \(a+b>0\), and
  coadjoint action
  \[
    z_1^p z_2^q\cdot\rho_{a,b}
    =
    -(pb-qa+p-q)\rho_{a-p+1,b-q+1},
  \]
  with negative-index targets equal to zero.  Among continuous
  total-degree-zero pairings, the residue pairing is determined up to one
  nonzero scalar, and this scalar is the only normalization ambiguity in
  this statement.

  \emph{Proof inputs.}
  Appendix~\ref{sec:app-matlis-principal-parts} centralizes the
  reduced Matlis calculation in the notation of this statement.
  Theorem~\ref{thm:canonical-residue-pairing} proves the residue
  pairing and its total-degree-zero scalar ambiguity.  Proposition~\ref{prop:principal-part-coadjoint}
  computes the explicit coadjoint action.  Theorem~\ref{thm:principal-part-coadjoint-uniqueness}
  gives the total-degree-zero Matlis-form uniqueness statement.  Theorem~\ref{thm:pbw-vs-deformation}
  and Theorem~\ref{thm:no-polynomial-realization-categorical} state the
  no-go result for replacing \(\mathcal P\) by polynomial descendants.
  Theorem~\ref{thm:app-matlis-rees-fourier-bridge} supplies the limited
  Fourier--Rees label bridge through the Rees lattice
  \(\bigoplus\hbar^{a+b}\C[[\hbar]]\rho_{a,b}\) and proves that it is not
  an \(\mathfrak h\)-module bridge.  Its flatness is the PBW normal-form
  flatness of the left cyclic \(\mathcal D_\hbar\)-modules
  \(\mathcal D_\hbar/(\partial_1,\partial_2)\) and
  \(\mathcal D_\hbar/(z_1,z_2)\).

  \emph{Scope and exclusions.} A same-action \(D_\hbar\), Fourier, or Rees
  comparison identifying the PBW descendant module with the Matlis
  coadjoint module is ruled out on the generic fibre by local nilpotence
  unless the Hamiltonian action or the boundary target is changed.  The
  Fourier-twisted label bridge does not supply such a same-action
  comparison.
\end{stmt}

\begin{stmt}[Statement: brane-line factorization currents]\label{thm:lane-factorization-current}
  \emph{Status.} \emph{proved degreewise stable}.  The proof is for
  compactly supported de Rham Hamiltonian currents and compactly
  supported distributional principal-part currents on intervals.  This
  statement is a strict support-local \(P_0\) current statement in the
  reduced Hamiltonian sector; local constancy is asserted only after the
  compact-support de Rham contraction of the Hamiltonian density sector.

  \emph{Exact hypotheses.} Intervals lie in the brane line \(\R\);
  \(\mathfrak h_I=\Omega^\bullet_c(I)\widehat\otimes\mathfrak h\);
  \(\mathcal P_I=\mathcal D'_c(I)\widehat\otimes\mathcal P\);
  extension by zero is the functoriality; and source coordinates in the
  displayed boundary operation are shifted-cotangent CE coordinates
  \(u_{a(t)dt\otimes f}\), not Hamiltonian ghost coordinates.  The CE
  source is cofactorization data: extension by zero on current Lie
  algebras induces pullback/coproduct on CE cochains, while covariant
  factorization belongs to the compactly supported current target.

  \emph{Local assertion.} For each interval \(I\subset\R\),
  \[
    \mathfrak h_I=\Omega^\bullet_c(I)\widehat\otimes\mathfrak h,\qquad
    \mathcal P_I=\mathcal D'_c(I)\widehat\otimes\mathcal P,\qquad
    \mathfrak g_I=\mathfrak h_I\ltimes\mathcal P_I[1].
  \]
  The Hamiltonian current sector admits an interval-wise locally constant
  \(P_0\) factorization equivalence.  On length-one shifted-cotangent
  CE coordinates the source is
  \[
    u_{a(t)dt\otimes f}\longmapsto
    B_f(a)=
    \int_I a(t)\,
    \overline{\operatorname{Tr}f(\phi_1(t),\phi_2(t))}\,dt .
  \]
  In the degree-zero notation used here \(a\in\Omega^0_c(I)\), the
  paired density is \(a(t)\,dt\), \(|u_{a(t)dt\otimes f}|=0\), and the
  de Rham-current pairing is
  \[
    \langle B\otimes\rho,\;a\otimes f\rangle
    =B(a)\operatorname{Res}(\rho f).
  \]
  For a regular current \(B=b(t)\,dt\) this is
  \(\int_I a(t)b(t)\,dt\,\operatorname{Res}(\rho f)\).  Orientation
  reversal contributes only the ordinary sign of the one-form density
  \(dt\); no extra sign is inserted in this statement.

  \emph{Proof inputs.}
  Appendix~\ref{sec:app-factorization-current-conventions} fixes the
  interval source, current dual, shifted-cotangent CE coordinate, and
  orientation convention used in this statement.
  Construction~\ref{constr:interval-fact-algebras} fixes the de
  Rham-current objects and extension-by-zero functoriality.
  Proposition~\ref{prop:brane-bracket-locality} proves disjoint-support
  locality.  Theorem~\ref{thm:hamiltonian-current-center-lift} supplies
  the smeared Hamiltonian central operation.  Theorem~\ref{thm:hamiltonian-current-factorization}
  identifies the interval-wise reduced Hamiltonian current
  central-operation map and its
  compatibility with factorization products.

  \emph{Scope and exclusions.} The strict current statement does not include the
  unreduced BV/factorization-centre morphism or its centrality homotopies
  against arbitrary open observables.
\end{stmt}

\begin{stmt}[Statement: reduced principal-part defect currents]\label{thm:lane-reduced-principal-part-current}
  \emph{Status.} \emph{proved degreewise stable}.  The proof applies
  only after the de Rham-current correction and only in the
  reduced \(P_0\) model after contraction.  It is not an unreduced
  BV/factorization theorem and does not supply centrality homotopies
  against arbitrary open observables.

  \emph{Exact hypotheses.} The algebra is the reduced
  defect-current model on intervals after de Rham-current contraction;
  the Hamiltonian factor is
  \(\widehat{\mathbf S}(\Omega_c^0(I)\widehat\otimes\bar A)\); the
  cotangent factor is
  \(\widehat{\mathbf S}(\mathcal D'_c(I)\widehat\otimes\mathcal P[1])\);
  brackets are strict \(P_0\)-brackets; the completed symmetric algebra
  is interpreted degreewise in the current category; and the primitive
  one-\(\psi\) comparison is taken only after passing to
  indecomposables.

  \emph{Local assertion.} The principal-part cotangent factor extends
  the Hamiltonian brane-line algebra to the reduced defect-current
  model
  \[
    A^{\mathrm{pp}}_\partial(I)
    =
    \widehat{\mathbf S}(\Omega_c^0(I)\widehat\otimes\bar A)
    \widehat\otimes
    \widehat{\mathbf S}(\mathcal D'_c(I)\widehat\otimes\mathcal P[1])
  \]
  with
  \[
    \{B_f(a),\Theta_\rho(B)\}
    =
    \Theta_{f\cdot\rho}(aB),
    \qquad
    \{\Theta_\rho(B),\Theta_\sigma(C)\}=0 .
  \]
  The chain-level primitive projection gives only the primitive
  indecomposable \(P_0\) shadow of this reduced comparison.
  Its quantum analogue is the reduced Moyal principal-part target of
  Construction~\ref{constr:quantum-principal-part-descendant-target},
  where \(f\cdot\rho\) is replaced by
  \(\operatorname{ad}^{\vee}_{f,\hbar}\rho\).

  \emph{Proof inputs.}
  Theorem~\ref{thm:reduced-principal-part-boundary-current} proves the
  reduced current model and its \(P_0\) bracket.  Theorem~\ref{thm:boundary-principal-part-cotangent-operators}
  records the boundary local-dual realization in the same reduced
  setting.  Remark~\ref{rmk:reduced-status-principal-part-currents} and
  Remark~\ref{rmk:unreduced-lift-status} state the unreduced
  obligations.  Theorem~\ref{thm:chain-level-primitive-projection},
  Corollary~\ref{cor:unreduced-d-closed}, and
  Remark~\ref{rmk:T5-primitive-shadow-status} supply only the primitive
  shadow status.  Proposition~\ref{prop:quantum-descendant-label-obstruction}
  and Theorem~\ref{thm:one-psi-label-hbar-obstruction} prove
  that the polynomial primitive label projection does not quantize to
  the Moyal coadjoint target.

  \emph{Scope and exclusions.} The unreduced cotangent factorization-centre
  morphism and its centrality homotopies remain outside this statement.  The
  reduced Moyal principal-part target is constructed, but a quantum lift
  from polynomial one-\(\psi\) primitives is ruled out under the
  leading-term assumptions above.  The pointwise classical obstruction
  is the mismatch between the polynomial one-\(\psi\) action and the
  principal-part coadjoint action:
  \[
    O_{1,0}:\widetilde\Psi_{a,b}\mapsto b\,\widetilde\Psi_{a,b-1},
    \qquad
    z_1:\rho_{a,b}\mapsto-(b+1)\rho_{a,b+1},
  \]
  and similarly
  \[
    O_{0,1}:\widetilde\Psi_{a,b}\mapsto
      -a\,\widetilde\Psi_{a-1,b},
    \qquad
    z_2:\rho_{a,b}\mapsto(a+1)\rho_{a+1,b}.
  \]
  Therefore the primitive projection records labels only; it does not
  supply centrality homotopies against arbitrary unreduced open
  observables.
\end{stmt}

\begin{stmt}[Statement: degree-zero Moyal/Weyl sector and graph boundary]\label{thm:lane-moyal-degree-zero}
  \emph{Status.} \emph{proved in finite algebraic model} for the formal
  Weyl/Moyal coefficient algebra and the reduced open-line Wick
  normalization.  The degree-zero stable trace comparison is
  \emph{conditional} on the external radial-parts input of
  Statement~\ref{stmt:app-radial-external-input}.  The mixed QME,
  unweighted regulator endpoint, descendant one-\(\psi\) lift, and
  Costello brane-defect graph realization are separate residual statements
  and are not hypotheses for the formal Moyal coefficient theorem.

  \emph{Exact hypotheses.} Work in the scalar-reduced degree-zero
  Hamiltonian sector at finite \(N\), assume the radial-parts input of
  Statement~\ref{stmt:app-radial-external-input} for the all-\(N\)
	  reduction, then take the stable connected trace range degreewise; use
	  the Capelli/Weyl normalization of the Weyl algebra and the formal
	  Moyal product on \(\bar A_\hbar\); and exclude one-antifield
	  descendants and graph-theoretic BV data from the proved comparison.
	  The all-\(N\) radial-parts input includes the matched-normalization
	  assertion \(J_N(f)=J_N^{\mathrm{rad}}(f)\); the formal Moyal
	  coefficient theorem and the cited radial-parts sources do not imply
	  this equality in the present Weyl/Capelli convention.

  \emph{Local assertion.} On the scalar-reduced degree-zero Hamiltonian
  sector, conditional on the radial-parts input, the classical
  closed-open comparison deforms to all orders in \(\hbar\).  The
  commutator is the formal Moyal commutator on \(\bar A_\hbar\) under
  the Capelli/Weyl normalization, and the stable connected trace map
  satisfies
  \[
    \hbar^{-1}[\partial_{\mathrm{bb},\hbar}(f),
      \partial_{\mathrm{bb},\hbar}(g)]_{\mathrm{conn,raw}}
    =
    \partial_{\mathrm{bb},\hbar}(\{f,g\}_\hbar).
  \]
  The first and third reduced open-line coefficients are
  \(\hbar P(f,g)\) and \(\hbar^3P^3(f,g)/24\).  A Costello-theoretic
  graph lift of those coefficients is conditional on the sector,
  propagator, counterterm, anomaly, and QME data named in the analytic
  graph problem.
  Here ``first'' and ``third'' refer to raw star-commutator powers
  \(\hbar^1\) and \(\hbar^3\).  After division by \(\hbar\), these become
  the normalized Hamiltonian-bracket coefficients in powers \(\hbar^0\)
  and \(\hbar^2\), namely \(P(f,g)\) and \(P^3(f,g)/24\).
  In the descendant direction the reduced target is not polynomial:
  Construction~\ref{constr:quantum-principal-part-descendant-target}
  gives the Moyal principal-part coadjoint complex, and
  Theorem~\ref{thm:one-psi-label-hbar-obstruction} rules out
  any strict \(\hbar\)-adic deformation of the classical
  \(\Psi_{a,b}\mapsto\rho_{a,b}\) projection satisfying its flatness,
  PBW-linear-action, Moyal-coadjoint-target, and leading-term hypotheses.

  \emph{Proof inputs.}
  Appendix~\ref{sec:app-radial-parts-moyal} separates the formal
  Moyal/Capelli calculation from the external radial-parts input.
  Theorem~\ref{thm:finite-n-reduced-moyal} uses that input to obtain
  the finite-\(N\) radial-parts identity and the stable degree-zero
  commutator.
  Corollary~\ref{cor:degree-zero-quantum-upgrade} and
  Theorem~\ref{thm:phi-hbar-all-order} assemble the stable
  Capelli-renormalized map.  Proposition~\ref{prop:moyal-monomial}
  proves the Moyal coefficient formula, with an independent finite exact
  arithmetic sweep as verification evidence.  Proposition~\ref{prop:conditional-boundary-product-normalization}
  and Theorem~\ref{thm:open-line-midpoint-graph-weights} compute
  the reduced open-line first and third coefficients.
  Theorem~\ref{thm:first-third-costello-normalizations} is conditional
  on Problem~\ref{prob:analytic-graph-realization}, and
  Problem~\ref{prob:weighted-rg-locality} carries the mixed
  brane-defect QME obstruction.  Theorem~\ref{thm:phi-hbar-trunc}
  gives the finite nilpotent obstruction to an all-order descendant
  lift.  Proposition~\ref{prop:quantum-descendant-label-obstruction},
  Construction~\ref{constr:quantum-principal-part-descendant-target},
  and Theorem~\ref{thm:one-psi-label-hbar-obstruction} isolate
  the one-\(\psi\) quantum target and the no-go for the polynomial
  primitive projection.

  \emph{Scope and exclusions.} Without the radial-parts source input the all-\(N\)
  trace identity is conditional; the finite exact check covers only finite sweeps,
  the rank-two Cartan shadow, and the direct \(N=2,3\) radial
  restrictions.  The all-order Costello graph realization, the mixed
  brane-defect QME, the unweighted regulator problem, and an unreduced
  one-antifield Moyal lift are not consequences of this algebraic statement.
  The reduced principal-part Moyal target is only a target-side
  construction.
\end{stmt}

\begin{stmt}[Statement: scalar \(U(1)\) anomaly]\label{thm:lane-u1-anomaly}
  \emph{Status.} \emph{proved in finite algebraic model}.  The statement
  records the scalar-reduction obstruction to lifting the reduced
  Hamiltonian comparison back to the unreduced constant direction.

  \emph{Exact hypotheses.} The closed algebra is the unreduced polynomial
  Hamiltonian algebra \(\mathfrak h_{\mathrm{poly}}=\C[z_1,z_2]\) and
  its quotient \(\bar A=\mathfrak h_{\mathrm{poly}}/\C\cdot1\); the open
  side is the finite-\(N\) trace algebra before and after quotienting
  \(\C\cdot\operatorname{Tr}(1)\); and the quantum comparison uses the
  same Capelli/Weyl normalization as the degree-zero Moyal theorem.

  \emph{Local assertion.} The scalar-reduction axis carries the
  distinguished class
  \[
    [\bar c]\in H^2_{\mathrm{Lie}}(\bar A,\C),
  \]
  represented by the constant Taylor coefficient of the Poisson
  bracket.  The same class appears on the open side as the
  centre-of-mass \(\mathfrak{gl}_N\) anomaly; quantum mechanically its
  finite-\(N\) trace-character representative is the Capelli contact
  \(\hbar N\).  All reduced Hamiltonian comparisons are made after
  killing constants; the unreduced scalar line is governed by this
  anomaly class.

  \emph{Proof inputs.}
  Lemma~\ref{lem:omega-cocycle} identifies the closed Poisson cocycle.
  Theorem~\ref{thm:u1-center-anomaly} proves the closed-side central
  extension statement.  Theorem~\ref{thm:u1-center-anomaly-open}
  computes the open \(U(1)\) centre.  Theorem~\ref{thm:quantum-classical-anomaly}
  identifies the classical cocycle with the Rees central extension whose
  finite-\(N\) trace-character value is the quantum Capelli shift.
  Corollary~\ref{cor:local-bulk-boundary-coupling} gives the explicit
  local central term in the smeared boundary bracket.

  \emph{Scope and exclusions.} The statement does not lift the reduced comparison to the
  unreduced scalar direction and does not imply compact Calabi--Yau,
  BKM, Igusa, or sister-volume anomaly statements.
\end{stmt}
